% 恒星光谱（综述）
% license CCBYSA3
% type Wiki

本文根据 CC-BY-SA 协议转载翻译自维基百科\href{https://en.wikipedia.org/wiki/Stellar_classification}{相关文章}。

在天文学中，恒星分类是根据恒星的光谱特性对恒星进行的分类。来自恒星的电磁辐射通过棱镜或衍射光栅被分解成光谱，呈现出类似彩虹的连续颜色，其间夹杂着光谱线。每一条光谱线都对应一种特定的化学元素或分子，而谱线的强度则反映该元素的丰度。不同光谱线强度的变化主要由恒星光球层的温度决定，尽管在某些情况下也确实存在真实的元素丰度差异。恒星的光谱型是一种简短的代码，主要概括了其电离状态，从而为光球层温度提供一个客观的度量。

目前，大多数恒星采用摩根–基南（Morgan–Keenan，MK）分类系统进行分类，该系统使用字母 O、B、A、F、G、K、M，按从最热（O 型）到最冷（M 型）的顺序排列。每一个字母类别又进一步用一个数字细分，其中 0 表示最热，9 表示最冷（例如 A8、A9、F0、F1 就构成了一个由热到冷的序列）。此外，该序列还扩展出 W、S、C 三个类别，用于那些不符合经典体系的恒星。一些恒星遗迹或质量明显偏离恒星范围的天体也被赋予了字母标记：D 表示白矮星，L、T、Y 表示棕矮星（以及部分系外行星）。

在 MK 系统中，还会在光谱型后加上一个用罗马数字表示的光度级。光度级依据恒星光谱中某些吸收线的宽度来判定，这些线宽随恒星大气密度变化，因此可用于区分巨星与矮星。其中，光度级 0 或 Ia+ 表示特超巨星，I 表示超巨星，II 表示亮巨星，III 表示普通巨星，IV 表示亚巨星，V 表示主序星，sd（或 VI）表示亚矮星，而 D（或 VII）表示白矮星。太阳的完整光谱分类为 G2V，这表明太阳是一颗主序星，其表面温度约为 5800 K。
\subsection{传统颜色描述}
传统的恒星颜色描述只考虑恒星光谱的峰值位置。然而实际上，恒星会在整个电磁波谱范围内辐射。由于所有光谱颜色叠加后呈现为白色，人眼实际看到的恒星颜色要比传统颜色描述所暗示的浅得多。这种“明亮度”的特性说明，将恒星简单地赋予某种光谱颜色可能具有误导性。在排除暗光条件下的颜色对比效应后，在一般观测条件中，并不存在真正的绿色、青色、靛色或紫色恒星。像太阳这样的“黄色”矮星实际上呈现为白色；“红矮星”呈现为较深的黄色或橙色；而“棕矮星”并不会真的呈现棕色，而是假想中在近距离观察时可能显得暗红色，或呈灰色/黑色。
